% ======================================================================================
% Nom     : reports/latex/sections/07_resultats.tex
% Rôle    : Section présentant les résultats expérimentaux.
% Auteur  : Maxime BRONNY
% Version : V1
% Cadre   : Réalisé dans le cadre du cours Techniques d'apprentissage artificiel
%           Master 1 Informatique Big Data
% Usage   : Fichier inclus dans main.tex lors de la compilation du rapport.
% ======================================================================================

\section{Résultats}
\label{sec:resultats}

% Toutes les valeurs de cette section proviennent de l'exécution réelle du
% pipeline (python3 main.py, 12/06/2026, split 70/15/15 stratifié, seed 42).
% Source : results/metrics.json et results/sklearn_comparison.json.

\subsection{Sélection des hyperparamètres sur la validation}

Le tableau~\ref{tab:validation} retrace la sélection effectuée sur l'ensemble
de validation (critère : F1-score). Pour l'arbre, le F1 progresse nettement de
la profondeur 3 à la profondeur 7, puis recule légèrement à 9 : au-delà de 7,
la complexité supplémentaire n'apporte plus de pouvoir prédictif généralisable.
Pour le $k$-NN, le F1 est remarquablement stable sur la grille, signe que la
décision dépend peu du voisinage exact au-delà d'une dizaine de voisins.

\begin{table}[H]
  \centering
  \caption{Sélection des hyperparamètres sur l'ensemble de validation
    (F1-score ; en gras, la configuration retenue).}
  \label{tab:validation}
  \begin{tabular}{llcc}
    \toprule
    \textbf{Modèle} & \textbf{Configuration} & \textbf{Seuil optimal} & \textbf{F1 (validation)} \\
    \midrule
    Régression logistique & lr $=0{,}1$, $1\,000$ it. & \textbf{0,3} & \textbf{0,625} \\
    \midrule
    \multirow{4}{*}{Arbre de décision}
      & profondeur 3 & 0,3 & 0,728 \\
      & profondeur 5 & 0,3 & 0,776 \\
      & \textbf{profondeur 7} & \textbf{0,6} & \textbf{0,819} \\
      & profondeur 9 & 0,6 & 0,817 \\
    \midrule
    \multirow{4}{*}{$k$-NN}
      & $k = 5$  & 0,5 & 0,683 \\
      & $k = 11$ & 0,4 & 0,695 \\
      & $k = 21$ & 0,3 & 0,692 \\
      & $\boldsymbol{k = 31}$ & \textbf{0,3} & \textbf{0,695} \\
    \bottomrule
  \end{tabular}
\end{table}

\subsection{Performances sur l'ensemble de test}

Le tableau~\ref{tab:resultats-test} présente l'évaluation finale, réalisée une
seule fois sur l'ensemble de test.

\begin{table}[H]
  \centering
  \caption{Métriques des trois modèles sur l'ensemble de test
    (4\,886 exemples). Hyperparamètres choisis sur la validation.}
  \label{tab:resultats-test}
  \begin{tabular}{lccccc}
    \toprule
    \textbf{Modèle} & \textbf{Accuracy} & \textbf{Précision} & \textbf{Rappel}
      & \textbf{F1-score} & \textbf{AUC-ROC} \\
    \midrule
    Régression logistique (seuil 0,3)        & 0,820 & 0,573 & 0,690 & 0,626 & 0,848 \\
    Arbre de décision (prof.\ 7, seuil 0,6)  & \textbf{0,931} & \textbf{0,977}
      & \textbf{0,703} & \textbf{0,817} & \textbf{0,909} \\
    $k$-NN ($k=31$, seuil 0,3)               & 0,862 & 0,683 & 0,684 & 0,683 & 0,877 \\
    \bottomrule
  \end{tabular}
\end{table}

La hiérarchie est nette et cohérente sur toutes les métriques : \textbf{l'arbre
de décision domine}, suivi du $k$-NN, puis de la régression logistique. L'arbre
atteint un F1 de $0{,}817$ et une AUC de $0{,}909$ ; son avance la plus
spectaculaire concerne la précision ($0{,}977$) : lorsqu'il prédit un défaut,
il se trompe très rarement. Les trois modèles plafonnent en revanche à un
rappel compris entre $0{,}69$ et $0{,}70$ : quel que soit le modèle, environ
$30\,\%$ des défauts échappent à la détection. Nous y revenons dans l'analyse
critique.

\subsection{Matrices de confusion}

\begin{figure}[H]
  \centering
  \includegraphics[width=0.32\textwidth]{lr_confusion_matrix.png}\hfill
  \includegraphics[width=0.32\textwidth]{dt_confusion_matrix.png}\hfill
  \includegraphics[width=0.32\textwidth]{knn_confusion_matrix.png}
  \caption{Matrices de confusion sur l'ensemble de test : régression
    logistique (gauche), arbre de décision (centre), $k$-NN (droite).}
  \label{fig:confusion-matrices}
\end{figure}

La lecture détaillée des matrices (figure~\ref{fig:confusion-matrices})
éclaire les différences de comportement. Sur les 4\,886 dossiers de test
(dont 1\,066 défauts) :

\begin{itemize}
  \item l'\textbf{arbre} ne commet que \textbf{18 faux positifs} (bons clients
    refusés) pour 749 défauts détectés, soit un profil très « prudent dans
    l'alerte » ; ses erreurs sont presque exclusivement des faux négatifs
    (317) ;
  \item la \textbf{régression logistique}, avec son seuil abaissé à $0{,}3$,
    détecte un nombre comparable de défauts (736) mais au prix de 549 faux
    positifs, soit trente fois plus que l'arbre ;
  \item le \textbf{$k$-NN} occupe une position intermédiaire (729 défauts
    détectés, 339 faux positifs).
\end{itemize}

Autrement dit, à rappel quasi égal, les trois modèles se distinguent presque
uniquement par leur taux de fausses alertes. C'est ce qui explique l'écart
de précision et donc de F1.

\subsection{Courbes ROC et comparaison globale}

\begin{figure}[H]
  \centering
  \includegraphics[width=0.75\textwidth]{roc_curves_comparison.png}
  \caption{Courbes ROC des trois modèles sur l'ensemble de test.
    AUC : arbre de décision $0{,}909$ ; $k$-NN $0{,}877$ ;
    régression logistique $0{,}848$.}
  \label{fig:roc}
\end{figure}

Les courbes ROC (figure~\ref{fig:roc}) confirment la hiérarchie
indépendamment du seuil de décision : la courbe de l'arbre domine celle du
$k$-NN, qui domine celle de la régression logistique, sur la quasi-totalité de
la plage de fonctionnement. Les trois modèles sont très au-dessus de la
diagonale du hasard (AUC $0{,}5$). La figure~\ref{fig:models-comparison}
récapitule les cinq métriques. L'ensemble des figures de ce rapport est produit
avec Matplotlib~\cite{hunter2007matplotlib}.

\begin{figure}[H]
  \centering
  \includegraphics[width=0.92\textwidth]{models_comparison.png}
  \caption{Comparaison des cinq métriques sur l'ensemble de test.}
  \label{fig:models-comparison}
\end{figure}

\subsection{Validation des implémentations avec scikit-learn}

Pour vérifier la justesse de nos implémentations, chaque modèle a été
ré-entraîné avec son équivalent scikit-learn sur exactement les mêmes ensembles
et avec des hyperparamètres équivalents (tableau~\ref{tab:sklearn}).

\begin{table}[H]
  \centering
  \caption{Vérification des implémentations from scratch : AUC sur l'ensemble
    de test, comparée à scikit-learn.}
  \label{tab:sklearn}
  \begin{tabular}{lccc}
    \toprule
    \textbf{Modèle} & \textbf{AUC from scratch} & \textbf{AUC scikit-learn}
      & \textbf{Écart} \\
    \midrule
    Régression logistique & 0,8479 & 0,8483 & $-0{,}0005$ \\
    Arbre de décision     & 0,9086 & 0,9086 & $0{,}0000$ \\
    $k$-NN                & 0,8772 & 0,8772 & $0{,}0000$ \\
    \bottomrule
  \end{tabular}
\end{table}

L'arbre et le $k$-NN obtiennent une AUC \emph{identique} à celle de
scikit-learn, et l'écart de la régression logistique ($5 \times 10^{-4}$)
s'explique par les différences d'optimiseur (notre descente de gradient pure
contre L-BFGS) et par la régularisation L2 activée par défaut dans
scikit-learn. Ces écarts négligeables valident nos trois implémentations.
L'ensemble du pipeline, comprenant le prétraitement, l'entraînement des trois
modèles, la sélection sur validation, l'évaluation et la génération des figures,
s'exécute en une vingtaine de secondes.
